%% 使用 jnuthesis 文档类生成南京大学学位论文的示例文档
%%
%% 作者：胡海星，starfish (at) gmail (dot) com
%% 项目主页: http://haixing-hu.github.io/jnu-thesis/
%%
%% 本样例文档中用到了吕琦同学的博士论文的提高和部分内容，在此对他表示感谢。
%%
\documentclass[master,winfonts]{jnuthesis}
%% jnuthesis 文档类的可选参数有：
%%   nobackinfo 取消封二页导师签名信息。注意，按照南大的规定，是需要签名页的。
%%   phd/master/bachelor 选择博士/硕士/学士论文

% 使用 blindtext 宏包自动生成章节文字
% 这仅仅是用于生成样例文档，正式论文中一般用不到该宏包
\usepackage[math]{blindtext}

%%%%%%%%%%%%%%%%%%%%%%%%%%%%%%%%%%%%%%%%%%%%%%%%%%%%%%%%%%%%%%%%%%%%%%%%%%%%%%%
% 设置论文的中文封面

% 论文标题，不可换行
% \title{}
% 如果论文标题过长，可以分两行，第一行用\titlea{}定义，第二行用\titleb{}定义，将上面的\title{}注释掉
\titlea{基于Docker, Android与Web的社会关系研究}
\titleb{——以江南大学为例}
% 论文的英文标题第一行
\englishtitlea{A Research on Social Relationships based on Docker, Android and Web}
% 论文的英文标题第二行
\englishtitleb{-- Take Jiangnan University as an Example}

% 论文作者姓名
\author{王雪瑞}
% 导师姓名职称
\supervisor{刘渊~~教授}
% 论文作者的学科与专业方向
\major{数字媒体技术}
% 论文作者的研究方向
\researchfield{网络媒体应用与安全技术}
% 论文作者所在院系的中文名称
\department{}
% 论文作者所在学校或机构的名称。此属性可选，默认值为``江南大学''。
\institute{江南大学}
% 论文的指导小组成员。
\directormembers{}
% 论文的提交日期，需设置年、月、日。
\submitdate{2017年5月}
% 论文的答辩日期，需设置年、月、日。
\defenddate{2017年6月}
% 论文的定稿日期，需设置年、月、日。此属性可选，默认值为最后一次编译时的日期，精确到日。
%% \date{2013年5月1日}
% 获得学位日期。
\degreedate{}

%%%%%%%%%%%%%%%%%%%%%%%%%%%%%%%%%%%%%%%%%%%%%%%%%%%%%%%%%%%%%%%%%%%%%%%%%%%%%%%
% 设置论文的英文封面

% 论文作者姓名的拼音
\englishauthor{WANG Xuerui}
% 导师姓名职称的英文
\englishsupervisor{}
% 论文作者学科与专业的英文名
\englishmajor{}
% 论文作者所在院系的英文名称
\englishdepartment{}
% 论文作者所在学校或机构的英文名称。此属性可选，默认值为``Jiangnan University''。
\englishinstitute{Jiangnan University}
% 论文完成日期的英文形式，它将出现在英文封面下方。需设置年、月、日。日期格式使用美国的日期
% 格式，即``Month day, year''，其中``Month''为月份的英文名全称，首字母大写；``day''为
% 该月中日期的阿拉伯数字表示；``year''为年份的四位阿拉伯数字表示。此属性可选，默认值为最后
% 一次编译时的日期。
\englishdate{May 1, 2017}

%%%%%%%%%%%%%%%%%%%%%%%%%%%%%%%%%%%%%%%%%%%%%%%%%%%%%%%%%%%%%%%%%%%%%%%%%%%%%%%
% 设置论文的中文摘要

% 设置中文摘要页面的论文标题及副标题的第一行。
% 此属性可选，其默认值为使用|\title|命令所设置的论文标题
% \abstracttitlea{数据中心网络模型研究}
% 设置中文摘要页面的论文标题及副标题的第二行。
% 此属性可选，其默认值为空白
% \abstracttitleb{}

%%%%%%%%%%%%%%%%%%%%%%%%%%%%%%%%%%%%%%%%%%%%%%%%%%%%%%%%%%%%%%%%%%%%%%%%%%%%%%%
% 设置论文的英文摘要

% 设置英文摘要页面的论文标题及副标题的第一行。
% 此属性可选，其默认值为使用|\englishtitle|命令所设置的论文标题
\englishabstracttitlea{}
% 设置英文摘要页面的论文标题及副标题的第二行。
% 此属性可选，其默认值为空白
\englishabstracttitleb{}

%%%%%%%%%%%%%%%%%%%%%%%%%%%%%%%%%%%%%%%%%%%%%%%%%%%%%%%%%%%%%%%%%%%%%%%%%%%%%%%
\begin{document}

%%%%%%%%%%%%%%%%%%%%%%%%%%%%%%%%%%%%%%%%%%%%%%%%%%%%%%%%%%%%%%%%%%%%%%%%%%%%%%%

% 制作中文封面
\maketitle

% 输出独创性声明与论文使用授权说明页面
\makeoriginalitypage

%%%%%%%%%%%%%%%%%%%%%%%%%%%%%%%%%%%%%%%%%%%%%%%%%%%%%%%%%%%%%%%%%%%%%%%%%%%%%%%
% 开始前言部分
\frontmatter

%%%%%%%%%%%%%%%%%%%%%%%%%%%%%%%%%%%%%%%%%%%%%%%%%%%%%%%%%%%%%%%%%%%%%%%%%%%%%%%
% 论文的中文摘要
\begin{abstract}
复杂网络的研究可上溯到20世纪60年代对ER网络的研究。90年后代随着Internet
的发展，以及对人类社会、通信网络、生物网络、社交网络等各领域研究的深入，
发现了小世界网络和无尺度现象等普适现象与方法。对复杂网络的定性定量的科
学理解和分析，已成为如今网络时代科学研究的一个重点课题。

在此背景下，由于云计算时代的到来，本文针对面向云计算的数据中心网络基础
设施设计中的若干问题，进行了几方面的研究。………………
% 中文关键词。关键词之间用中文全角分号隔开，末尾无标点符号。
\keywords{小世界理论；网络模型；数据中心}
\end{abstract}

%%%%%%%%%%%%%%%%%%%%%%%%%%%%%%%%%%%%%%%%%%%%%%%%%%%%%%%%%%%%%%%%%%%%%%%%%%%%%%%
% 论文的英文摘要
\begin{englishabstract}
\blindtext
% 英文关键词。关键词之间用英文半角逗号隔开，末尾无符号。
\englishkeywords{Small World, Network Model, Data Center}
\end{englishabstract}

%%%%%%%%%%%%%%%%%%%%%%%%%%%%%%%%%%%%%%%%%%%%%%%%%%%%%%%%%%%%%%%%%%%%%%%%%%%%%%%
% 生成论文目次
\tableofcontents

%%%%%%%%%%%%%%%%%%%%%%%%%%%%%%%%%%%%%%%%%%%%%%%%%%%%%%%%%%%%%%%%%%%%%%%%%%%%%%%
% 开始正文部分
\mainmatter

%%%%%%%%%%%%%%%%%%%%%%%%%%%%%%%%%%%%%%%%%%%%%%%%%%%%%%%%%%%%%%%%%%%%%%%%%%%%%%%
% 学位论文的正文应以《绪论》作为第一章
\chapter{绪论}\label{chapter_introduction}
\section{研究背景}

在分布式网络领域，沿着高性能集群、普世计算、网格计算的方向，现已走入云
计算时代。

云计算对信息技术架构造成了越来越大的影响。例如，借助Amazon EC2云平台，
用户借助其基础设施，可以十分方便的部署各类应用，以支持企业服务需求。用
户可以按需购买计算资源，网络带宽，存储空间等各类资源以支持他们的业务需
求，并在业务完成之后迅速的归还这些资源。通过云技术，用户可以集中在他们
擅长的核心业务之中，而不会被诸如硬件购买、安装系统、网络设置、备份和安
全等等问题干扰。

与此同时，随着计算机的普及化和微型化，现在的手持设备拥有不输于7年前台式
机的处理能力。在网络时代面前，智能终端广泛普及，每个人都可成为信息源。
在信息爆炸的时代，数据挖掘、机器学习、金融分析和模拟等行业中不断涌现新
的需求，诸如针对用户行为和社会关系的挖掘进行广告精准投放，用户行为预测
等。为了支撑PB级尺度的数据规模，需要海量的计算节点，催生并不断促进了各
行各业对云计算基础设施的建设需求。

海量的数据需要海量的处理能力，然而海量的处理能力又需要高带宽的网络IO为
承载。作为云环境中最基础的一环，IaaS层在网络、存储、计算资源的分割这几
方面，承担起整个系统的基石。虽然并非必须，但一般来说，为考虑沙盘环境，
以及对资源的细粒度切割分配，IaaS通常会伴随着虚拟化技术的运用。虚拟化具
有许多与云计算切合的特点，例如，虚拟化可以屏蔽物理环境的差异，可在多物
理节点中进行无缝迁移，可对系统进行快照和还原。这些特点都与云计算时代所
追求的灵活性、高伸缩性、快速响应等特点而吻合。

在虚拟化实现方面，目前已取得了诸如ESXi，Xen，KVM等成熟成果。然而当虚拟
化扩大的一定规模，随着节点数目的增多，在网络方面将会面临一系列取舍的问
题。例如，基于二层交换的扁平网络，当节点数目上升到千数量级时，广播报文
将会极大的拖累网络性能，必须通过划分子网，通过三层路由等形式重新规划为
多层网络结构；另一方面，除了联通之外，还需要考虑ACL控制，负载均衡，外网
通讯等各类防火墙以及NAT规则的实现。这些复杂的网络配置，在一定程度上抵消
了虚拟机带来的灵活性。例如虽然虚拟机可根据需要动态迁移，但在迁移之后，
由于网络位置的变化需要重新进行网络参数配置。虽然虚拟机在迁移过程中系统
内部状态没有变化，但站在网络角度看，该虚拟节点跟关机重启没有区别。

针对上述问题，本文站在面向云计算时代的数据中心网络建设的角度，对网络模
型进行深入研究和探讨。通过改善二层交换网络的ARP机制来解决广播风暴问题，
引入比树形网络更为复杂的复杂网络理论，指导网络节点的互联模型。从而将网
络的复杂性隐藏在节点环境之外，在节点层面仅提供简单但巨大的二层交换扁平
网络。

\section{研究目的与意义}
\subsection{现有问题与不足}

测试一下引用\cite{newman2006structure}，连续引用
\cite{newman2001random,aiello2000random,bollobas2001random}，另一个连续引用
\cite{newman2001random,bollobas2001random,barabasi1999emergence}。测试一下带页码
的引用\cite[124--128]{erdHos1961strength}。

\subsection{中心观点与思想}

云计算在概念上通常被分为IaaS、PaaS、SaaS几个层面。但透过分类去理解其本
质，可认为是上世纪70年代基于大型计算机的中心控制型瘦客户端终端模式，在
如今技术水平上的一种新的表达，是在技术发展道路中，螺旋上升的结果。

与瘦客户端相比，云计算在设计结构上存在一定的相似性。

在设计思路上，两者都为了降低管理成本和硬件成本、以低能耗、高弹性等需求
为设计目标。随着技术的进步，云计算在具体实现形态上与传统的大型机也有很
大的不同：

一方面，云中心不再是传统的一台大型机，而是用大量廉价计算节点的互联来提
供海量资源。云计算更强调资源规模的无缝、平滑扩展，以及高可靠性，无单点
故障问题。另一方面，云计算时代的终端，也具备相当计算能力。随着web2.0的
整合，还有向胖客户端和智能终端发展的趋势。

总而言之，云计算在大框架中是传统的中心控制／终端的模式，但在中心与终端
两方面，都引入分布式技术加以改良。核心的思路是在低成本的前提下做到高可
靠性、高灵活性和高伸缩性。因此，云计算并不仅仅以数量换性能的表象，本质
上为低成本高性能，追求高能效比，并在实现层面讲究可实现性和可操作性。

%%%%%%%%%%%%%%%%%%%%%%%%%%%%%%%%%%%%%%%%%%%%%%%%%%%%%%%%%%%%%%%%%%%%%%%%%%%%%%%
% 学位论文的正文应以《结论》作为最后一章
\chapter{结论}\label{chapter_concludes}

本文在第\ref{chapter_smallworld}章中，通过考虑数据中心网络布局构建中的最大度限制
问题，提出了符合数据中心网络基本要求的DS小世界模型，并分析了它的性质。随后提出
SIDN，将DS模型映射到具体的网络结构中，并分析了所构成网络的平均直径、网络总带宽、
对故障的容错能力等各项网络性能。

分析与仿真实验证明，SIDN网络具有很好的扩展能力，网络总带宽与网络规模成
近似线性增长的关系；具有很强的容错能力，链路损坏与节点损坏几乎无法破坏
网络的联通性，故障率对网络性能的影响与破坏节点/链路占总资源比率线性相关。

随后在第\ref{chapter_scalefree}章中，分析了无尺度网络在数据中心网络构建应用中的
理论方面问题。对Scafida \cite{gyarmati2010scafida}文中所述在最大度限制的情况下运
用BA算法构造的网络并不会损失无尺度性质的观点，进行了深入的分析，并指出了该论点的
局限性。

在给出了在引入节点最大度限制之后，利用分治和递归的思想，对无尺度网络
进行多层构建，对所构造的网络进行度-度相关性，以及聚类性分析。

\begin{table}
  \centering
  \begin{tabular}{cccp{38mm}}
    \toprule
    \textbf{文档域类型} & \textbf{Java类型} & \textbf{宽度(字节)} & \textbf{说明} \\
    \midrule
    BOOLEAN  & boolean &  1  & \\
    CHAR     & char    &  2  & UTF-16字符 \\
    BYTE     & byte    &  1  & 有符号8位整数 \\
    SHORT    & short   &  2  & 有符号16位整数 \\
    INT      & int     &  4  & 有符号32位整数 \\
    LONG     & long    &  8  & 有符号64位整数 \\
    STRING   & String  &  字符串长度  & 以UTF-8编码存储 \\
    DATE     & java.util.Date & 8 & 距离GMT时间1970年1月1日0点0分0秒的毫秒数 \\
    BYTE\_ARRAY & byte$[]$ & 数组长度 & 用于存储二进制值 \\
    BIG\_INTEGER & java.math.BigInteger & 和具体值有关 & 任意精度的长整数 \\
    BIG\_DECIMAL & java.math.BigDecimal & 和具体值有关 & 任意精度的十进制实数 \\
    \bottomrule
  \end{tabular}
  \caption{测试表格}\label{table:test5}
\end{table}

表\ref{table:test5}用于测试表格。随后分析了无尺度网络构造过程中，交换机节点与数
据节点的角色区别，分析了两者在不同比率下形成的网络形态，以及对网络性能造成的影响。

通过理论分析和仿真实验，分析并找出比率因子q的最佳取值。此外，无尺度现象
的引入提高了网络的聚类系数，从而在不失灵活性可靠性的基础上，进一步提升
了网络的性能。

在第\ref{chapter_random}章中，将关注点转移到交换机本身。由于图论难以描述数据中心
网络中的交换设备，因此放弃基于图的抽象模型，转而基于多维簇划分的思想，提出并设计
了WarpNet网络模型。

该网络模型突破了基于图描述的局限性，并对网络的带宽等指标进行理论分析并
给出定量描述。最后对比了理论分析、仿真测试结果，并在实际物理环境中进系
真实部署，通过6节点的小规模实验以及1000节点虚拟机的大规模实验，表明该模
型的理论分析、仿真测试与实际实验吻合，并在网络性能、容错能力、伸缩性灵
活性方面得到了进一步的提升。

在第\ref{chapter_experiments}章中，针对网络模型研究这一类工作的共性，设计构造通
用验证平台系统。以海量虚拟机和虚拟分布式交换机的形式，实现了基于少量物理节点，对
大规模节点的模拟。其模拟运行的过程与真实运行在实现层面完全一致，运行的结果与真实
环境线性相关。除为本文所涉若干网络模型提供验证外，可进一步推广到更为广泛的领域，
为各种网络模型及路由算法的研究工作，提供分析、指导与验证。

%%%%%%%%%%%%%%%%%%%%%%%%%%%%%%%%%%%%%%%%%%%%%%%%%%%%%%%%%%%%%%%%%%%%%%%%%%%%%%%
% 致谢，应放在《结论》之后
\begin{acknowledgement}
  首先感谢我的母亲韦春花对我的支持。其次感谢我的导师陈近南对我的精心指导和热心帮助。接下来，
  感谢我的师兄茅十八和风际中，他们阅读了我的论文草稿并提出了很有价值的修改建议。

  最后，感谢我亲爱的老婆们：双儿、苏荃、阿珂、沐剑屏、曾柔、建宁公主、方怡，感谢
  你们在生活上对我无微不至的关怀和照顾。我爱你们！
\end{acknowledgement}

%%%%%%%%%%%%%%%%%%%%%%%%%%%%%%%%%%%%%%%%%%%%%%%%%%%%%%%%%%%%%%%%%%%%%%%%%%%%%%%
% 附录
\appendix

\chapter{博士(硕士)学位论文编写格式规定(试行)}

\section{适用范围}

本规定适用于博士学位论文编写，硕士学位论文编写应参照执行。

\section{引用标准}

GB7713科学技术报告、学位论文和学术论文的编写格式。

GB7714文后参考文献著录规则。

\section{印制要求}

论文必须用白色纸印刷，并用A4(210mm×297mm)标准大小的白纸。纸的四周应留足空白
边缘，上方和左侧应空边25mm以上，下方和右侧应空边20mm以上。除前置部分外，其它
部分双面印刷。

论文装订不要用铁钉，以便长期存档和收藏。

论文封面与封底之间的中缝（书脊）必须有论文题目、作者和学校名。

\section{编写格式}

论文由前置部分、主体部分、附录部分(必要时)、结尾部分(必要时)组成。

前置部分包括封面，题名页，声明及说明，前言，摘要(中、英文)，关键词，目次页，
插图和附表清单(必要时)，符号、标志、缩略词、首字母缩写、单位、术语、名词解释
表(必要时)。

主体部分包括绪论(作为正文第一章)、正文、结论、致谢、参考文献表。

附录部分包括必要的各种附录。

结尾部分包括索引和封底。

\section{前置部分}

\subsection{封面（博士论文国图版用）}

封面是论文的外表面，提供应有的信息，并起保护作用。

封面上应包括下列内容：
\begin{enumerate}
\item 分类号  在左上角注明分类号，便于信息交换和处理。一般应注明《中国图书资
  料分类法》的类号，同时应注明《国际十进分类法UDC》的类号；
\item 密级  在右上角注明密级；
\item “博士学位论文”用大号字标明；
\item 题名和副题名   用大号字标明；
\item 作者姓名；
\item 学科专业名称；
\item 研究方向；
\item 导师姓名，职称；
\item 日期包括论文提交日期和答辩日期；
\item 学位授予单位。
\end{enumerate}

\subsection{题名}

题名是以最恰当、最简明的词语反映论文中最重要的特定内容的逻辑组合。

题名所用每一词语必须考虑到有助于选定关键词和编写题录、索引等二次文献可以提供
检索的特定实用信息。

题名应避免使用不常见的缩略词、首字母缩写字、字符、代号和公式等。

题名一般不宜超过20字。

论文应有外文题名，外文题名一般不宜超过10个实词。

可以有副题名。

题名在整本论文中不同地方出现时，应完全相同。

\subsection{前言}

前言是作者对本论文基本特征的简介，如论文背景、主旨、目的、意义等并简述本论文
的创新性成果。

\subsection{摘要}

摘要是论文内容不加注释和评论的简单陈述。

论文应有中、英文摘要，中、英文摘要内容应相同。

摘要应具有独立性和自含性，即不阅读论文的全文，便能获得必要的信息，摘要
中有数据、有结论，是一篇完整的短文，可以独立使用，可以引用，可以用于推广。摘
要的内容应包括与论文同等量的主要信息，供读者确定有无必要阅读全文，也供文摘等
二次文献引用。摘要的重点是成果和结论。

中文摘要一般在1500字，英文摘要不宜超过1500实词。

摘要中不用图、表、化学结构式、非公知公用的符号和术语。

\subsection{关键词}

关键词是为了文献标引工作从论文中选取出来用于表示全文主题内容信息款目的单词或
术语。

每篇论文选取3－8个词作为关键词，以显著的字符另起一行，排在摘要的左下方。在英
文摘要的左下方应标注与中文对应的英文关键词。

\subsection{目次页}

目次页由论文的章、节、附录等的序号、名称和页码组成，另页排在摘要的后面。

\subsection{插图和附表清单}

论文中如图表较多，可以分别列出清单并置于目次页之后。

图的清单应有序号、图题和页码。表的清单应有序号、表题和页码。

符号、标志、缩略词、首字母缩写、计量单位、名词、术语等的注释表符号、标志、缩略词、
首字母缩写、计量单位、名词、术语等的注释说明汇集表，应置于图表清单之后。

\section{主体部分}

\subsection{格式}

主体部分由绪论开始，以结论结束。主体部分必须由另页右页开始。每一章必须另页开
始。全部论文章、节、目的格式和版面安排要划一，层次清楚。

\subsection{序号}

\begin{figure}[htbp]
  \centering
  \includegraphics[width= 0.5\textwidth]{jnuname.eps}\\
  \caption{测试附录中的插图}\label{fig:appendix1}
\end{figure}

论文的章可以写成：第一章。节及节以下均用阿拉伯数字编排序号，如
1.1，1.1.1等。

论文中的图、表、附注、参考文献、公式、算式等一律用阿拉伯数字分别分章依序连续编排
序号。其标注形式应便于互相区别，一般用下例：图1.2；表2.3；附注1）；文献[4]；式
  (6.3)等。

论文一律用阿拉伯数字连续编页码。页码由首页开始，作为第1页，并为右页另页。封页、
封二、封三和封底不编入页码，应为题名页、前言、目次页等前置部分单独编排页码。页码
必须标注在每页的相同位置，便于识别。

\begin{equation}
    C_i = \frac{2E_i}{k_i(k_i-1)}
\end{equation}

附录依序用大写正体A、B、C、$\cdots$编序号，如：附录A。附录中的图、表、式、参考文
献等另行编序号，与正文分开，也一律用阿拉伯数字编码，但在数码前题以附条序码，如图
A.1；表B.2；式(B.3)；文献[A.5]等。

\subsection{绪论}

绪论（综述）：简要说明研究工作的目的、范围、相关领域的前人工作和知识空白、理
论基础和分析，研究设想、研究方法和实验设计、预期结果和意义等。一般在教科书中
有的知识，在绪论中不必赘述。

绪论的内容应包括论文研究方向相关领域的最新进展、对有关进展和问题的评价、本论
文研究的命题和技术路线等；绪论应表明博士生对研究方向相关的学科领域有系统深入
的了解，论文具有先进性和前沿性；

\begin{problem}
测试定理环境。测试定理环境。测试定理环境。测试定理环境。测试定理环境。测试定理环境。
测试定理环境。测试定理环境。测试定理环境。
\end{problem}

为了反映出作者确已掌握了坚实的基础理论和系统的专门知识，具有开阔的科学视野，对研
究方案作了充分论证，绪论应单独成章，列为第一章，绪论的篇幅应达$1\sim 2$万字，不
得少于$1$万字；绪论引用的文献应在$100$篇以上，其中外文文献不少于$60\%$；引用文献
应按正文中引用的先后排列。

\subsection{正文}

论文的正文是核心部分，占主要篇幅。正文必须实事求是，客观真切，准确完备，合乎
逻辑，层次分明，简便可读。

\begin{figure}[htbp]
  \centering
  \includegraphics[width= 0.5\textwidth]{jnuname.eps}\\
  \caption{测试附录中的插图}\label{fig:appendix2}
\end{figure}

正文的每一章(除绪论外)应有小结，在小结中应明确阐明作者在本章中所做的工作，特
别是创新性成果。凡本论文要用的基础性内容或他人的成果不应单独成章，也不应作过
多的阐述，一般只引结论、使用条件等，不作推导。

\subsubsection{图}

图包括曲线图、构造图、示意图、图解、框图、流程图、记录图、布置图、地图、照片
、图版等。

图应具有“自明性”，即只看图、图题和图例，不阅读正文，就可以理解图意。

图应编排序号。每一图应有简短确切的图题，连同图号置于图下。必要时，应将图上的
符号、标记、代码，以及实验条件等，用最简练的文字，横排于图题下方，作为图例说
明。

\begin{example}
测试定理环境。测试定理环境。测试定理环境。测试定理环境。测试定理环境。测试定理环境。
测试定理环境。测试定理环境。测试定理环境。
\end{example}

曲线图的纵、横坐标必须标注“量、标准规定符号、单位”。此三者只有在不必要标明
(如无量纲等)的情况下方可省略。坐标上标注的量的符号和缩略词必须与正文一致。

照片图要求主题和主要显示部分的轮廓鲜明，便于制版。如用放大缩小的复制品，必须
清晰，反差适中。照片上应该有表示目的物尺寸的标度。

\subsubsection{表}

表的编排，一般是内容和测试项目由左至右横读，数据依序竖排。表应有自明性。

表应编排序号。

每一表应有简短确切的表题，连同标号置于表上。必要时，应将表中的符号、标记、代
码，以及需要说明事项，以最简练的文字，横排于表题下，作为表注，也可以附注于表
下。表内附注的序号宜用小号阿拉伯数字并加圆括号置于被标注对象的右上角，如：
xxx${}^{1)}$；不宜用“*”，以免与数学上共轭和物质转移的符号相混。

表的各栏均应标明“量或测试项目、标准规定符号、单位”。只有在无必要标注的情况下
方可省略。表中的缩略词和符号，必须与正文中一致。

表内同一栏的数字必须上下对齐。表内不宜用“同上”，“同左”和类似词，一律填入具体数字
或文字。表内“空白”代表未测或无此项，“－”或“\textellipsis”（因“－”可能与代表阴性
  反应相混）代表未发现，“0”代表实测结果确为零。

如数据已绘成曲线图，可不再列表。

\subsubsection{数学、物理和化学式}

正文中的公式、算式或方程式等应编排序号，序号标注于该式所在行(当有续行时，应
标注于最后一行)的最右边。

较长的式，另行居中横排。如式必须转行时，只能在$+$，$-$，$\times$，$\div$，$<$，
$>$处转行。上下式尽可能在等号“$=$”处对齐。

小数点用“$.$”表示。大于$999$的整数和多于三位数的小数，一律用半个阿拉伯数字符的小
间隔分开，不用千位撇。对于纯小数应将$0$列于小数点之前。

示例：应该写成$94\ 652.023\ 567$和$0.314\ 325$, 不应写成$94,652.023,567$和
$.314,325$。

应注意区别各种字符，如：拉丁文、希腊文、俄文、德文花体、草体；罗马数字和阿拉伯数
字；字符的正斜体、黑白体、大小写、上下脚标（特别是多层次，如“三踏步”）、上下偏差
等。

\subsubsection{计量单位}

报告、论文必须采用国务院发布的《中华人民共和国法定计量单位》，并遵照《中华人
民共和国法定计量单位使用方法》执行。使用各种量、单位和符号，必须遵循附录B所
列国家标准的规定执行。单位名称和符号的书写方式一律采用国际通用符号。

\subsubsection{符号和缩略词}

符号和缩略词应遵照国家标准的有关规定执行。如无标准可循，可采纳本学科或本专业
的权威性机构或学术团体所公布的规定；也可以采用全国自然科学名词审定委员会编印
的各学科词汇的用词。如不得不引用某些不是公知公用的、且又不易为同行读者所理解
的、或系作者自定的符号、记号、缩略词、首字母缩写字等时，均应在第一次出现时一
一加以说明，给以明确的定义。

\subsection{结论}

报告、论文的结论是最终的、总体的结论，不是正文中各段的小结的简单重复。结论应
该准确、完整、明确、精炼。在结论中要清楚地阐明论文中有那些自己完成的成果，特
别是创新性成果；

如果不可能导出应有的结论，也可以没有结论而进行必要的讨论。可以在结论或讨论中
提出建议、研究设想、仪器设备改进意见、尚待解决的问题等。

\subsection{致谢}

可以在正文后对下列方面致谢：

\begin{itemize}
\item 国家科学基金、资助研究工作的奖学金基金、合作单位、资助或支持的企业、组织或个
人；
\item 协助完成研究工作和提供便利条件的组织或个人；
\item 在研究工作中提出建议和提供帮助的人；
\item 给予转载和引用权的资料、图片、文献、研究思想和设想的所有者；
\item 其他应感谢的组织或个人。
\end{itemize}

\subsection{参考文献表}

\subsubsection{专著著录格式}

主要责任者，其他责任者，书名，版本，出版地：出版者，出版年

例：1. 刘少奇，论共产党员的修养，修订2版，北京：人民出版社，1962

\subsubsection{连续出版物中析出的文献著录格式}

析出文献责任者，析出文献其他责任者，析出题名，原文献题名，版本：文献中的位置。

例：2. 李四光，地壳构造与地壳运动，中国科学，1973 (4)：400－429

参考文献采用顺序编码制，按论文正文所引用文献出现的先后顺序连续编码。

\section{附录}

附录是作为报告、论文主体的补充项目，并不是必需的。

下列内容可以作为附录编于报告、论文后，也可以另编成册；

\begin{enumerate}
\item 为了整篇论文材料的完整，但编入正文又有损于编排的条理和逻辑性，这一材料
包括比正文更为详尽的信息、研究方法和技术更深入的叙述，建议可以阅读的参考文献
题录，对了解正文内容有用的补充信息等；
\item 由于篇幅过大或取材于复制品而不便于编入正文的材料；
\item 不便于编入正文的罕见珍贵资料；
\item 对一般读者并非必要阅读，但对本专业同行有参考价值的资料；
\item 某些重要的原始数据、数学推导、计算程序、框图、结构图、注释、统计表、计
算机打印输出件等。
\end{enumerate}

附录与正文连续编页码。

每一附录均另页起。

\section{结尾部分 (必要时)}

为了将论文迅速存储入电子计算机，可以提供有关的输入数据。可以编排分类索引、著者索
引、关键词索引等。

% 参考文献。应放在\backmatter之前。
% 推荐使用BibTeX，若不使用BibTeX时注释掉下面一句。
\nocite{*}
\bibliography{sample}
% 不使用 BibTeX
%\begin{thebibliography}{2}
%
%\bibitem{deng:01a}
%{邓建松,彭冉冉,陈长松}.
%\newblock {\em \LaTeXe{}科技排版指南}.
%\newblock 科学出版社,书号:7-03-009239-2/TP.1516, 北京, 2001.
%
%\bibitem{wang:00a}
%王磊.
%\newblock {\em \LaTeXe{}插图指南}.
%\newblock 2000.
%\end{thebibliography}

% 附录，必须放在参考文献后，backmatter前
\appendix
\chapter{图论基础知识}
\Blindtext

%%%%%%%%%%%%%%%%%%%%%%%%%%%%%%%%%%%%%%%%%%%%%%%%%%%%%%%%%%%%%%%%%%%%%%%%%%%%%%%
% 书籍附件
\backmatter
%%%%%%%%%%%%%%%%%%%%%%%%%%%%%%%%%%%%%%%%%%%%%%%%%%%%%%%%%%%%%%%%%%%%%%%%%%%%%%%
% 作者简历与科研成果页，应放在backmatter之后
\begin{resume}
% 论文作者身份简介，一句话即可。
\begin{authorinfo}
\noindent 韦小宝，男，汉族，1985年11月出生，江苏省扬州人。
\end{authorinfo}
% 论文作者教育经历列表，按日期从近到远排列，不包括将要申请的学位。
\begin{education}
\item[2007年9月 --- 2010年6月] 南京大学计算机科学与技术系 \hfill 硕士
\item[2003年9月 --- 2007年6月] 南京大学计算机科学与技术系 \hfill 本科
\end{education}
% 论文作者在攻读学位期间所发表的文章的列表，按发表日期从近到远排列。
\begin{publications}
\item Xiaobao Wei, Jinnan Chen, ``Voting-on-Grid Clustering for Secure
  Localization in Wireless Sensor Networks,'' in \textsl{Proc. IEEE International
    Conference on Communications (ICC) 2010}, May. 2010.
\item Xiaobao Wei, Shiba Mao, Jinnan Chen, ``Protecting Source Location Privacy
  in Wireless Sensor Networks with Data Aggregation,'' in \textsl{Proc. 6th
    International Conference on Ubiquitous Intelligence and Computing (UIC)
    2009}, Oct. 2009.
\end{publications}
% 论文作者在攻读学位期间参与的科研课题的列表，按照日期从近到远排列。
\begin{projects}
\item 国家自然科学基金面上项目``无线传感器网络在知识获取过程中的若干安全问题研究''
（课题年限~2010年1月 --- 2012年12月），负责位置相关安全问题的研究。
\item 江苏省知识创新工程重要方向项目下属课题``下一代移动通信安全机制研究''
（课题年限~2010年1月 --- 2010年12月），负责LTE/SAE认证相关的安全问题研究。
\end{projects}
\end{resume}

%%%%%%%%%%%%%%%%%%%%%%%%%%%%%%%%%%%%%%%%%%%%%%%%%%%%%%%%%%%%%%%%%%%%%%%%%%%%%%%
% 生成《学位论文出版授权书》页面，应放在最后一页
\makelicense

%%%%%%%%%%%%%%%%%%%%%%%%%%%%%%%%%%%%%%%%%%%%%%%%%%%%%%%%%%%%%%%%%%%%%%%%%%%%%%%
\end{document}
